\chapter{Complexity — Tautology}
%%% generated by the blueprint pipeline from TCSlib.Complexity.ClassNP.Tautology
%%%%%%%%%%%%%%%%%%%%%%%%%%%%%%%%%%%%%%%%%%%%%%%%%%%%%%%%%%%%%%%%%%%%%%%%%%%%%%%
\section{Overview}

This module defines coNP-hardness and coNP-completeness
(Karp-reduction form) and the language $\mathrm{TAUTOLOGY}$ on the DNF fragment, and
proves [AB09, Example~2.21]: $\mathrm{TAUTOLOGY}$ is coNP-complete.
Membership is witnessed by a concrete linear-time verifier machine that validates the
formula syntax and then evaluates the DNF against a falsifying-assignment
certificate; hardness composes the Cook--Levin reduction to $\mathrm{SAT}$ with a
streaming transducer that takes the De Morgan dual of the decoded CNF by flipping
every literal polarity in place.

Throughout this chapter a formula over $\mathbb{N}$-indexed variables is a list of
lists of literals $(v,b)$ with variable index $v \in \mathbb{N}$ and polarity
$b \in \{0,1\}$; a literal is satisfied by an assignment $a : \mathbb{N} \to \{0,1\}$
when $a(v) = b$. Read as a CNF, the formula is an AND of ORs (clauses); read as a
DNF, it is an OR of ANDs (terms), so the empty DNF is false and an empty term is
true. The serialization of a formula writes each clause or term as a leading $1$,
its serialized literals — a literal $(v,b)$ becomes $1^{v+1}\,0\,b$ — and a closing
$0$, with one final $0$ terminating the formula. Decoding is total: a string that is
not exactly the serialization of some formula decodes to the empty formula. The
variable bound of a formula is one more than its largest variable index ($0$ if no
literal occurs). The pair encoding $\mathrm{pair}(a,u)$ of two binary words doubles
every bit of $a$, then appends the separator $01$ followed by $u$. For a binary word
$u$, the certificate assignment $\alpha_u$ maps $v$ to the $v$-th bit of $u$ when
$v < |u|$ and to $0$ otherwise.

%%%%%%%%%%%%%%%%%%%%%%%%%%%%%%%%%%%%%%%%%%%%%%%%%%%%%%%%%%%%%%%%%%%%%%%%%%%%%%%
\section{Declarations}

\begin{definition}[coNP-hardness]
\lean{Complexity.coNPHard}
\label{Complexity.coNPHard}
\leanok
\uses{Complexity.coNP}
A language $L$ over the binary alphabet is coNP-hard if every language
$L'$ in coNP is polynomial-time Karp reducible to $L$, that is,
$L' \le_p L$ [AB09, Section~2.6.1].
\end{definition}

\begin{definition}[coNP-completeness]
\lean{Complexity.coNPComplete}
\label{Complexity.coNPComplete}
\leanok
\uses{Complexity.coNP, Complexity.coNPHard}
A language $L$ over the binary alphabet is coNP-complete if $L$ belongs
to coNP and $L$ is coNP-hard, i.e.\ every language in
coNP is polynomial-time Karp reducible to $L$ [AB09, Section~2.6.1].
\end{definition}

\begin{definition}[The language TAUTOLOGY]
\lean{Complexity.TAUTOLOGY}
\label{Complexity.TAUTOLOGY}
\leanok
\uses{Std.Sat.DNF.Tautology, Std.Sat.DNF.decode}
The language $\mathrm{TAUTOLOGY}$ is the set of binary strings $x$ whose decoded
formula, read as a DNF (an OR of ANDs of literals), evaluates to true under every
assignment $a : \mathbb{N} \to \{0,1\}$ [AB09, Example~2.21, DNF fragment]. Since
decoding is total with the empty formula as fallback, and the empty disjunction is
false, no non-well-formed string belongs to $\mathrm{TAUTOLOGY}$.
\end{definition}

\begin{lemma}[Polarity flip preserves the variable bound]
\lean{Complexity.taut_numVars_dual}
\label{Complexity.taut_numVars_dual}
\leanok
\uses{Std.Sat.CNF.numVars, Std.Sat.DNF, Std.Sat.DNF.dual}
\difficulty{2}
For every DNF formula $\psi$, the CNF obtained from $\psi$ by flipping the polarity of
every literal has the same variable bound as the term list of $\psi$ viewed as a CNF,
where the variable bound of a formula is one more than its largest variable index
(and $0$ if no literal occurs).
\end{lemma}

\begin{lemma}[DNF evaluation depends only on mentioned variables]
\lean{Complexity.taut_eval_congr}
\label{Complexity.taut_eval_congr}
\leanok
\uses{Complexity.eval_congr_of_lt_numVars, Complexity.taut_numVars_dual, Std.Sat.CNF.eval_dual, Std.Sat.CNF.numVars, Std.Sat.DNF, Std.Sat.DNF.dual, Std.Sat.DNF.dual_dual, Std.Sat.DNF.eval}
\difficulty{3}
Let $\psi$ be a DNF formula and let $a, b : \mathbb{N} \to \{0,1\}$ be assignments
with $a(v) = b(v)$ for every $v$ below the variable bound of $\psi$ (one more than its
largest variable index). Then $\psi$ evaluates to the same truth value under $a$ and
under $b$.
\end{lemma}

\begin{definition}[Assignment carried by a certificate]
\lean{Complexity.tautAssignment}
\label{Complexity.tautAssignment}
\leanok
Given a binary word $u$, the certificate assignment $\alpha_u : \mathbb{N} \to \{0,1\}$
maps a variable index $v$ to the $v$-th bit of $u$ when $v < |u|$, and to $0$ (false)
when $v \ge |u|$.
\end{definition}

\begin{lemma}[Certificate characterization of non-tautologies]
\lean{Complexity.taut_certificate_equiv}
\label{Complexity.taut_certificate_equiv}
\leanok
\uses{Complexity.TAUTOLOGY, Complexity.tautAssignment, Complexity.taut_eval_congr, Std.Sat.CNF.numVars_decode_le, Std.Sat.DNF.decode, Std.Sat.DNF.eval}
\difficulty{4}
For every binary string $x$, $x$ lies outside $\mathrm{TAUTOLOGY}$ if and only if
there exists a binary word $u$ of length exactly $|x| + 1$ such that the DNF decoded
from $x$ evaluates to false under the assignment $\alpha_u$ that reads variable $v$
off the $v$-th bit of $u$ (false for $v \ge |u|$).
\end{lemma}

\begin{lemma}[Split success equation]
\lean{Complexity.taut_split_some}
\label{Complexity.taut_split_some}
\leanok
\uses{Turing.solveSplit}
\difficulty{2}
Let $N$ and $i$ be natural numbers. If the bounded split search with parameters
$(1,1)$ — which returns the least $i \le N$ with $i + (i+1) = N$, if any — returns
$i$ on $N$, then $i + (i + 1) = N$.
\end{lemma}

\begin{lemma}[Split completeness]
\lean{Complexity.taut_split_exists}
\label{Complexity.taut_split_exists}
\leanok
\uses{Complexity.taut_split_some, Turing.solveSplit}
\difficulty{3}
If natural numbers $i$ and $N$ satisfy $i + (i + 1) = N$, then the bounded split
search with parameters $(1,1)$, which looks for the least $j \le N$ with
$j + (j+1) = N$, succeeds on $N$ and returns exactly $i$.
\end{lemma}

\begin{definition}[Verifier bit for non-tautologies]
\lean{Complexity.tautVerifierBit}
\label{Complexity.tautVerifierBit}
\leanok
\uses{Complexity.tautAssignment, Std.Sat.DNF.decode, Std.Sat.DNF.eval, Turing.solveSplit}
Given a binary string $z$, the verifier bit is false when $|z|$ cannot be written as
$i + (i + 1)$; otherwise, for the unique such $i$, it is the negation of the truth
value of the DNF decoded from the first $i$ bits of $z$ under the assignment that
reads variable $v$ off the $v$-th bit of the remaining suffix of $z$ (false beyond
its end).
\end{definition}

\begin{definition}[Verifier language for non-tautologies]
\lean{Complexity.tautVerifier}
\label{Complexity.tautVerifier}
\leanok
\uses{Complexity.tautVerifierBit}
The verifier language is the set of binary strings $z$ such that $|z| = i + (i + 1)$
for some $i$ and the DNF decoded from the first $i$ bits of $z$ evaluates to false
under the assignment reading variable $v$ off the $v$-th bit of the remaining $i + 1$
bits (false beyond them).
\end{definition}

\begin{lemma}[Verifier on a correctly sized concatenation]
\lean{Complexity.taut_verifier_append}
\label{Complexity.taut_verifier_append}
\leanok
\uses{Complexity.tautAssignment, Complexity.tautVerifier, Complexity.tautVerifierBit, Complexity.taut_split_exists, Std.Sat.DNF.decode, Std.Sat.DNF.eval}
\difficulty{3}
Let $x$ and $u$ be binary words with $|u| = |x| + 1$, and let $V$ be the set of binary
strings $z$ with $|z| = i + (i+1)$ whose first $i$ bits decode to a DNF that is false
under the assignment read off the remaining bits. Then
$x \mathbin{+\!\!+} u \in V$ if and only if the DNF decoded from $x$ evaluates to false
under the assignment $\alpha_u$ reading variable $v$ off the $v$-th bit of $u$ (false
beyond $|u|$).
\end{lemma}

\begin{lemma}[A polynomial-time verifier gives coNP membership]
\lean{Complexity.taut_membership_of_verifier}
\label{Complexity.taut_membership_of_verifier}
\leanok
\uses{Complexity.P, Complexity.TAUTOLOGY, Complexity.coNP, Complexity.tautVerifier, Complexity.taut_certificate_equiv, Complexity.taut_verifier_append}
\difficulty{3}
Let $V$ be the set of binary strings $z$ with $|z| = i + (i + 1)$ for some $i$ such
that the DNF decoded from the first $i$ bits of $z$ is false under the assignment
reading variable $v$ off the $v$-th of the remaining bits (false beyond them). If $V$
belongs to $\mathrm{P}$, then $\mathrm{TAUTOLOGY}$ belongs to coNP.
\end{lemma}

\begin{definition}[Syntax scanner states]
\lean{Complexity.TautSyntax}
\label{Complexity.TautSyntax}
\leanok
The control states of the finite-state syntax scanner for the formula serialization
grammar are six: formula level, clause level, unary index, polarity, exact end, and
error.
\end{definition}

\begin{definition}[Syntax scanner transition]
\lean{Complexity.tautSyntaxStep}
\label{Complexity.tautSyntaxStep}
\leanok
\uses{Complexity.TautSyntax}
The transition function $\delta$ of the six-state syntax scanner for the formula
serialization grammar is the following: at formula level, $1$ goes to clause level and
$0$ to exact end; at clause level, $1$ goes to unary index and $0$ back to formula
level; in the unary index, $1$ stays and $0$ goes to polarity; from polarity, any bit
returns to clause level; from exact end or error, any bit goes to error.
\end{definition}

\begin{definition}[Syntax scan acceptance]
\lean{Complexity.tautSyntaxAccept}
\label{Complexity.tautSyntaxAccept}
\leanok
\uses{Complexity.TautSyntax, Complexity.tautSyntaxStep}
Given a scanner state $q$ of the six-state formula syntax scanner and a binary string
$x$, the acceptance bit is true exactly when folding the scanner transition over all
of $x$, starting from $q$, ends in the exact-end state.
\end{definition}

\begin{lemma}[One scanner transition]
\lean{Complexity.taut_syntax_cons}
\label{Complexity.taut_syntax_cons}
\leanok
\uses{Complexity.TautSyntax, Complexity.tautSyntaxAccept, Complexity.tautSyntaxStep}
\difficulty{1}
For every state $q$ of the six-state formula syntax scanner with transition $\delta$,
every bit $b$, and every binary string $x$, the acceptance bit of the scan from $q$ on
$b :: x$ equals the acceptance bit of the scan from $\delta(q, b)$ on $x$.
\end{lemma}

\begin{lemma}[The error state never recovers]
\lean{Complexity.taut_syntax_bad}
\label{Complexity.taut_syntax_bad}
\leanok
\uses{Complexity.tautSyntaxAccept, Complexity.tautSyntaxStep, Complexity.taut_syntax_cons}
\difficulty{3}
For every binary string $x$, the acceptance bit of the six-state formula syntax
scanner started in its error state and run over $x$ is false.
\end{lemma}

\begin{lemma}[Exact end accepts only the empty string]
\lean{Complexity.taut_syntax_done}
\label{Complexity.taut_syntax_done}
\leanok
\uses{Complexity.tautSyntaxAccept, Complexity.taut_syntax_bad}
\difficulty{2}
For every binary string $x$, the six-state formula syntax scanner started in its
exact-end state accepts $x$ if and only if $x$ is empty.
\end{lemma}

\begin{lemma}[Unary scanner partition]
\lean{Complexity.taut_takeTrues_shape}
\label{Complexity.taut_takeTrues_shape}
\leanok
\uses{Std.Sat.CNF.takeTrues}
\difficulty{3}
For every binary string $x$, if the leading-ones scanner returns the count $k$ and the
suffix $s$ on $x$, then $x = 1^{k} \mathbin{+\!\!+} s$.
\end{lemma}

\begin{lemma}[Literal parsing reconstructs its input]
\lean{Complexity.taut_parseLit_shape}
\label{Complexity.taut_parseLit_shape}
\leanok
\uses{Complexity.taut_takeTrues_shape, Std.Sat.CNF.parseLit, Std.Sat.CNF.serializeLit}
\difficulty{3}
If the literal parser succeeds on a binary string $x$, returning a literal $\ell$ and
a remainder $r$, then $x$ is exactly the serialization of $\ell$ followed by $r$.
\end{lemma}

\begin{lemma}[Clause parsing reconstructs its input]
\lean{Complexity.taut_parseClause_shape}
\label{Complexity.taut_parseClause_shape}
\leanok
\uses{Complexity.taut_parseLit_shape, Std.Sat.CNF.parseClause, Std.Sat.CNF.parseLit, Std.Sat.CNF.serializeClause}
\difficulty{5}
For every fuel value, if the fuel-indexed clause parser succeeds on a binary string
$x$, returning a clause $C$ and a remainder $r$, then $x$ is exactly the serialization
of $C$ (including its closing $0$) followed by $r$.
\end{lemma}

\begin{lemma}[Formula parsing reconstructs its input]
\lean{Complexity.taut_parseClauses_shape}
\label{Complexity.taut_parseClauses_shape}
\leanok
\uses{Complexity.taut_parseClause_shape, Std.Sat.CNF.parseClause, Std.Sat.CNF.parseClauses, Std.Sat.CNF.serialize}
\difficulty{5}
For every fuel value, if the fuel-indexed formula parser succeeds on a binary string
$x$, returning a formula $\varphi$ and a remainder $r$, then $x$ is exactly the
serialization of $\varphi$ followed by $r$.
\end{lemma}

\begin{lemma}[Exact parsing is inverse serialization]
\lean{Complexity.taut_parse_shape}
\label{Complexity.taut_parse_shape}
\leanok
\uses{Complexity.taut_parseClauses_shape, Std.Sat.CNF.parse, Std.Sat.CNF.serialize}
\difficulty{3}
If the exact-consumption formula parser succeeds on a binary string $x$, returning
the formula $\varphi$, then $x$ equals the serialization of $\varphi$.
\end{lemma}

\begin{lemma}[Syntax scan of a unary index]
\lean{Complexity.taut_syntax_unary}
\label{Complexity.taut_syntax_unary}
\leanok
\uses{Complexity.tautSyntaxAccept, Complexity.tautSyntaxStep, Complexity.taut_syntax_cons, Std.Sat.CNF.takeTrues}
\difficulty{3}
For every binary string $x$, write $x = 1^k s$ where $s$ is the suffix left by the
leading-ones scanner. The six-state formula syntax scan from the unary-index state
accepts $x$ if and only if $s = 0\,b\,r$ for some bit $b$ and string $r$ and the scan
from the clause-level state accepts $r$; if $s$ has fewer than two bits or starts
with $1$, the scan rejects.
\end{lemma}

\begin{lemma}[Syntax scan agrees with the literal parser]
\lean{Complexity.taut_syntax_literal}
\label{Complexity.taut_syntax_literal}
\leanok
\uses{Complexity.tautSyntaxAccept, Complexity.tautSyntaxStep, Complexity.taut_syntax_cons, Complexity.taut_syntax_unary, Std.Sat.CNF.parseLit, Std.Sat.CNF.takeTrues}
\difficulty{4}
For every binary string $x$, the string $1 :: x$ is accepted by the six-state formula
syntax scan from the clause-level state if and only if the literal parser succeeds on
$1 :: x$ with some remainder $r$ and $r$ is accepted by the scan from the clause-level
state; when the literal parser fails, the scan rejects.
\end{lemma}

\begin{lemma}[Syntax scan agrees with the clause parser]
\lean{Complexity.taut_syntax_clause}
\label{Complexity.taut_syntax_clause}
\leanok
\uses{Complexity.tautSyntaxAccept, Complexity.taut_parseLit_shape, Complexity.taut_syntax_literal, Std.Sat.CNF.parseClause, Std.Sat.CNF.parseLit, Std.Sat.CNF.serializeLit}
\difficulty{5}
Let $f$ be a fuel value and $x$ a binary string with $|x| \le f$. The six-state formula
syntax scan from the clause-level state accepts $x$ if and only if the clause parser
with fuel $f$ succeeds on $x$ with some remainder $r$ and the scan from the
formula-level state accepts $r$; when the clause parser fails, the scan rejects.
\end{lemma}

\begin{lemma}[Syntax scan agrees with the formula parser]
\lean{Complexity.taut_syntax_formula}
\label{Complexity.taut_syntax_formula}
\leanok
\uses{Complexity.tautSyntaxAccept, Complexity.tautSyntaxStep, Complexity.taut_parseClause_shape, Complexity.taut_syntax_clause, Complexity.taut_syntax_cons, Complexity.taut_syntax_done, Std.Sat.CNF.parseClause, Std.Sat.CNF.parseClauses}
\difficulty{5}
Let $f$ be a fuel value and $x$ a binary string with $|x| \le f$. The six-state formula
syntax scan from the formula-level state accepts $x$ if and only if the formula parser
with fuel $f$ succeeds on $x$ with an empty remainder.
\end{lemma}

\begin{lemma}[Syntax scan decides parseability]
\lean{Complexity.taut_syntax_parse}
\label{Complexity.taut_syntax_parse}
\leanok
\uses{Complexity.tautSyntaxAccept, Complexity.taut_syntax_formula, Std.Sat.CNF.parse, Std.Sat.CNF.parseClauses}
\difficulty{3}
For every binary string $x$, the string $x$ is accepted by the six-state formula
syntax scan from the formula-level state if and only if the exact-consumption formula
parser succeeds on $x$.
\end{lemma}

\begin{definition}[Evaluation states of the verifier]
\lean{Complexity.TautEval}
\label{Complexity.TautEval}
\leanok
The evaluation states of the non-tautology verifier are formula level, and the states
clause level, unary index, and polarity, which each carry a Boolean accumulator
recording whether every literal of the current term read so far is satisfied.
\end{definition}

\begin{definition}[Control states of the verifier]
\lean{Complexity.TautControl}
\label{Complexity.TautControl}
\leanok
\uses{Complexity.TautEval, Complexity.TautSyntax}
The control states of the non-tautology verifier machine are: syntax states carrying a
formula-scanner state, either before the first copy of a doubled input bit or holding
the first copy already read; certificate-copy and copy-back states; two input-rewind
states; evaluation states carrying an evaluation state, either before a doubled bit or
holding its first copy; a certificate-rewind state carrying the term accumulator; and
a verdict state carrying the output bit.
\end{definition}

\begin{definition}[Silent verifier move]
\lean{Complexity.tautMove}
\label{Complexity.tautMove}
\leanok
\uses{Complexity.TautControl, Turing.Action}
Given an input-head displacement $d$ and a work-head displacement $e$ in
$\{-1, 0, 1\}$ and a control state $q$, the silent move is the one-work-tape machine
action that moves the input head by $d$ and the work head by $e$, writes nothing,
emits nothing, and enters $q$.
\end{definition}

\begin{definition}[Evaluation action on a formula bit]
\lean{Complexity.tautEvalAction}
\label{Complexity.tautEvalAction}
\leanok
\uses{Complexity.TautControl, Complexity.TautEval, Complexity.tautMove, Turing.Action}
Given an evaluation state, a formula bit $b$, and the work cell $w$ under the
certificate head (blank read as $0$), the evaluation action is a silent move that
advances the input head and acts as follows. At formula level, $1$ opens a term with
accumulator true and $0$ enters the accepting verdict. At clause level with
accumulator $c$, $1$ starts a literal and $0$ ends the term, entering the rejecting
verdict if $c$ is true and returning to formula level otherwise. In the unary index,
each $1$ advances the certificate head and $0$ moves to polarity; the polarity bit
$b$ sets the accumulator to $c \wedge (w = b)$, moves the certificate head left, and
enters the certificate-rewind state.
\end{definition}

\begin{definition}[The non-tautology verifier machine]
\lean{Complexity.tautTM}
\label{Complexity.tautTM}
\leanok
\uses{Complexity.TautControl, Complexity.tautEvalAction, Complexity.tautMove, Complexity.tautSyntaxStep, Turing.FinTM}
The non-tautology verifier is a finite Turing machine over the binary alphabet with
one work tape. On an input $\mathrm{pair}(x, u)$ it reads the doubled bits of $x$ two
at a time through the six-state formula syntax scanner; at the separator it emits $1$
and halts if $x$ is not well formed, and otherwise copies $u$ onto the work tape,
rewinds both heads, and evaluates the DNF serialized by $x$ term by term against the
certificate assignment on the work tape, finally emitting the negation of its truth
value. An input ending before a separator, or a doubled pair $10$, makes it emit $0$;
no bit is emitted before the final verdict.
\end{definition}

\begin{definition}[Indexed verifier configuration]
\lean{Complexity.tautCfg}
\label{Complexity.tautCfg}
\leanok
\uses{Complexity.TautControl, Turing.Cfg, Turing.FinTM.bufferTape}
Given an input $x$, an optional control state $q$ (absent meaning halted), an index
$i \le |x|$, a word $u$, an integer $h$, and an output word (empty by default), the
indexed configuration is the configuration of the one-work-tape non-tautology
verifier on $x$ with state $q$, the input head having consumed $i$ input bits, the
work tape holding $u$ from cell $0$ (blank elsewhere), the work head at cell $h$,
and the given output.
\end{definition}

\begin{lemma}[Input symbol at an indexed configuration]
\lean{Complexity.taut_read}
\label{Complexity.taut_read}
\leanok
\uses{Complexity.TautControl, Complexity.tautCfg, Turing.Cfg.inputSymbol, Turing.FinTM.inputSymbol_at}
\difficulty{1}
Let $x$ be an input and $i \le |x|$. In every configuration of the non-tautology
verifier on $x$ whose input head has consumed $i$ bits, the symbol under the input
head is the $i$-th bit of $x$ when $i < |x|$ and the blank when $i = |x|$,
regardless of state, work tape, and output.
\end{lemma}

\begin{lemma}[Silent right move]
\lean{Complexity.taut_move_right}
\label{Complexity.taut_move_right}
\leanok
\uses{Complexity.TautControl, Complexity.tautCfg, Complexity.tautMove, Turing.Action.apply, Turing.moveInputPos_pos_of_ne_right}
\difficulty{3}
Let $x$ be an input with $i + 1 \le |x|$. Applying the silent action that moves the
input head right, moves the work head by $e$, and enters state $q'$ to a
non-tautology verifier configuration on $x$ in any state with $i$ input bits consumed, work word
$u$, work head $h$, and output $o$ yields the configuration in state $q'$ with $i+1$
bits consumed, the same $u$ and $o$, and work head $h + e$.
\end{lemma}

\begin{lemma}[Silent stationary move]
\lean{Complexity.taut_move_stay}
\label{Complexity.taut_move_stay}
\leanok
\uses{Complexity.TautControl, Complexity.tautCfg, Complexity.tautMove, Turing.Action.apply}
\difficulty{2}
Let $x$ be an input and $i \le |x|$. Applying the silent action that keeps the input
head fixed, moves the work head by $e$, and enters state $q'$ to a non-tautology
verifier configuration on $x$ in any state with $i$ input bits consumed, work word $u$, work head
$h$, and output $o$ yields the configuration in state $q'$ with the same index, $u$,
and $o$, and work head $h + e$.
\end{lemma}

\begin{lemma}[A doubled bit is one syntax step]
\lean{Complexity.taut_syntax_double}
\label{Complexity.taut_syntax_double}
\leanok
\uses{Complexity.TautSyntax, Complexity.tautCfg, Complexity.tautSyntaxStep, Complexity.tautTM, Complexity.taut_move_right, Complexity.taut_read, Turing.Action.apply, Turing.MultiTapeTM.runFrom, Turing.MultiTapeTM.runFrom_succ_eq_step, Turing.MultiTapeTM.runFrom_zero, Turing.MultiTapeTM.step}
\difficulty{4}
Let $x$ be an input with $i + 2 \le |x|$ whose bits at positions $i$ and $i + 1$ both
equal $b$, and let $s$ be a formula-scanner state with transition $\delta$. Started in
the syntax state carrying $s$ with $i$ input bits consumed, empty work tape, work head
$0$, and empty output, the non-tautology verifier reaches, after exactly two steps,
the syntax state carrying $\delta(s, b)$ with $i + 2$ bits consumed and all else
unchanged.
\end{lemma}

\begin{lemma}[The pair separator ends the syntax pass]
\lean{Complexity.taut_syntax_separator}
\label{Complexity.taut_syntax_separator}
\leanok
\uses{Complexity.TautSyntax, Complexity.tautCfg, Complexity.tautTM, Complexity.taut_move_right, Complexity.taut_read, Turing.Action.apply, Turing.MultiTapeTM.runFrom, Turing.MultiTapeTM.runFrom_succ_eq_step, Turing.MultiTapeTM.runFrom_zero, Turing.MultiTapeTM.step}
\difficulty{4}
Let $x$ be an input with $i + 2 \le |x|$ whose bits at positions $i$ and $i+1$ are $0$
and $1$, and let $s$ be a formula-scanner state. Started in the syntax state carrying
$s$ with $i$ input bits consumed, empty work tape, head $0$, and empty output, after
exactly two steps the non-tautology verifier has consumed $i+2$ bits and is in the
certificate-copy state if $s$ is the exact-end state, and in the verdict state with
bit $1$ otherwise.
\end{lemma}

\begin{lemma}[The complete syntax pass]
\lean{Complexity.taut_syntax_run}
\label{Complexity.taut_syntax_run}
\leanok
\uses{Complexity.TautSyntax, Complexity.tautCfg, Complexity.tautSyntaxStep, Complexity.tautTM, Complexity.taut_syntax_double, Complexity.taut_syntax_separator, Turing.MultiTapeTM.runFrom, Turing.MultiTapeTM.runFrom_add, Turing.pairEncode, Turing.universal_pair_length}
\difficulty{5}
Let $a, u, p$ be binary words, $x = p \mathbin{+\!\!+} \mathrm{pair}(a, u)$, and $s$ a
formula-scanner state. Started in the syntax state carrying $s$ with $|p|$ input bits
consumed, empty work tape, head $0$, and empty output, the non-tautology verifier
after exactly $2|a| + 2$ steps has consumed $|p| + 2|a| + 2$ bits, with work tape and
output still empty, and is in the certificate-copy state if folding the scanner
transition over $a$ from $s$ ends in the exact-end state, and in the verdict state
with bit $1$ otherwise.
\end{lemma}

\begin{lemma}[Initial configuration of the verifier]
\lean{Complexity.taut_init}
\label{Complexity.taut_init}
\leanok
\uses{Complexity.tautCfg, Complexity.tautTM, Turing.Cfg.init, Turing.MultiTapeTM.initCfg}
\difficulty{2}
For every input $x$, the initial configuration of the non-tautology verifier on $x$ is
the configuration in the syntax state carrying the formula-level scanner state, with
no input bits consumed, an empty work tape, work head at cell $0$, and empty output.
\end{lemma}

\begin{lemma}[The verdict transition]
\lean{Complexity.taut_verdict}
\label{Complexity.taut_verdict}
\leanok
\uses{Complexity.tautCfg, Complexity.tautTM, Turing.Action.apply, Turing.MultiTapeTM.step}
\difficulty{2}
Let $x$ be an input, $i \le |x|$, $u$ a word, $h$ an integer, and $b$ a bit. From the
non-tautology verifier's configuration in the verdict state with bit $b$, $i$ input
bits consumed, work word $u$, work head $h$, and empty output, one step halts with
output exactly the single bit $b$, all other components unchanged.
\end{lemma}

\begin{lemma}[Malformed formulas are accepted]
\lean{Complexity.taut_machine_malformed}
\label{Complexity.taut_machine_malformed}
\leanok
\uses{Complexity.tautSyntaxAccept, Complexity.tautSyntaxStep, Complexity.tautTM, Complexity.taut_init, Complexity.taut_syntax_parse, Complexity.taut_syntax_run, Complexity.taut_verdict, Std.Sat.CNF.parse, Turing.FinTM.ComputesInTime, Turing.FinTM.computesInTime_iff, Turing.MultiTapeTM.runFrom_succ_eq_step', Turing.pairEncode}
\difficulty{4}
Let $x$ be a binary string on which the exact-consumption formula parser fails, and
let $u$ be any binary word. Then the non-tautology verifier, run on
$\mathrm{pair}(x, u)$, halts within $2|x| + 3$ steps with output the single bit $1$.
\end{lemma}

\begin{lemma}[The empty input is rejected]
\lean{Complexity.taut_machine_empty}
\label{Complexity.taut_machine_empty}
\leanok
\uses{Complexity.tautTM, Turing.FinTM.ComputesInTime, Turing.FinTM.computesInTime_iff}
\difficulty{2}
The non-tautology verifier, run on the empty input, halts within $2$ steps, and its
output is the single bit $0$.
\end{lemma}

\begin{lemma}[Certificate copy run]
\lean{Complexity.taut_copy_run}
\label{Complexity.taut_copy_run}
\leanok
\uses{Complexity.TautControl, Complexity.tautCfg, Complexity.tautTM, Complexity.taut_read, Turing.Action, Turing.Action.apply, Turing.FinTM.bufferTape_append, Turing.MultiTapeTM.runFrom, Turing.MultiTapeTM.runFrom_succ_eq_step', Turing.moveInputPos_pos_of_ne_right}
\difficulty{5}
Let the input be $x = p \mathbin{+\!\!+} u$. Started in the certificate-copy state with
$|p|$ input bits consumed, empty work tape, work head $0$, and empty output, for every
$j \le |u|$ the non-tautology verifier after exactly $j$ steps is still in the copy
state with $|p| + j$ bits consumed, the first $j$ bits of $u$ on the work tape, the
work head at cell $j$, and empty output.
\end{lemma}

\begin{lemma}[Work-tape rewind]
\lean{Complexity.taut_work_rewind}
\label{Complexity.taut_work_rewind}
\leanok
\uses{Complexity.TautControl, Complexity.tautCfg, Complexity.tautMove, Complexity.tautTM, Complexity.taut_move_stay, Turing.Action.apply, Turing.Cfg.workTapeSymbols, Turing.MultiTapeTM.runFrom, Turing.MultiTapeTM.runFrom_succ_eq_step, Turing.MultiTapeTM.runFrom_zero, Turing.MultiTapeTM.step}
\difficulty{5}
Let $q$ and $d$ be control states of the non-tautology verifier such that, in state
$q$, it silently moves the work head left with the input head fixed whenever the work
cell is nonblank, and otherwise moves the work head right and enters $d$. Then for
every input $x$, index $i \le |x|$, work word $u$, and $n \le |u|$, starting in state
$q$ with work head at cell $n - 1$, exactly $n + 1$ steps reach state $d$ with work
head at cell $0$, all other components unchanged.
\end{lemma}

\begin{lemma}[Startup of the evaluation phase]
\lean{Complexity.taut_evaluation_start}
\label{Complexity.taut_evaluation_start}
\leanok
\uses{Complexity.tautCfg, Complexity.tautSyntaxStep, Complexity.tautTM, Complexity.taut_copy_run, Complexity.taut_init, Complexity.taut_move_stay, Complexity.taut_read, Complexity.taut_syntax_parse, Complexity.taut_syntax_run, Complexity.taut_work_rewind, Std.Sat.CNF.parse, Turing.Action.apply, Turing.FinTM.timed_rewind, Turing.MultiTapeTM.initCfg, Turing.MultiTapeTM.runFrom, Turing.MultiTapeTM.runFrom_add, Turing.MultiTapeTM.runFrom_succ_eq_step', Turing.MultiTapeTM.step, Turing.pairEncode, Turing.universal_pair_length}
\difficulty{6}
Let $x$ be a binary string on which the exact-consumption formula parser succeeds, and
let $u$ be any binary word. Then for some $t \le 4(|\mathrm{pair}(x,u)| + 2)$, running
the non-tautology verifier for $t$ steps from its initial configuration on
$\mathrm{pair}(x,u)$ reaches the formula-level evaluation state with no input bits
consumed, $u$ on the work tape from cell $0$, the work head at cell $0$, and empty
output.
\end{lemma}

\begin{lemma}[A doubled bit is one evaluation step]
\lean{Complexity.taut_eval_double}
\label{Complexity.taut_eval_double}
\leanok
\uses{Complexity.TautControl, Complexity.TautEval, Complexity.tautCfg, Complexity.tautEvalAction, Complexity.tautMove, Complexity.tautTM, Complexity.taut_move_right, Complexity.taut_read, Turing.Action.apply, Turing.FinTM.bufferTape, Turing.MultiTapeTM.runFrom, Turing.MultiTapeTM.runFrom_succ_eq_step, Turing.MultiTapeTM.runFrom_zero, Turing.MultiTapeTM.step}
\difficulty{4}
Let $x$ be an input with $i + 2 \le |x|$ and $i$-th bit $b$, let $u$ be a word, $h$ an
integer, and $q$ an evaluation state. Suppose the evaluation action for $q$ on bit $b$,
with the work cell at $h$ of the tape holding $u$, is the silent move advancing the
input head, moving the work head by $e$, and entering $q'$. Then from the
non-tautology verifier's evaluation configuration carrying $q$ with $i$ bits consumed,
work word $u$, and head $h$, two steps reach state $q'$ with $i+2$ bits consumed and
head $h + e$.
\end{lemma}

\begin{lemma}[Pair encoding of a concatenation]
\lean{Complexity.taut_pair_append}
\label{Complexity.taut_pair_append}
\leanok
\uses{Turing.pairEncode}
\difficulty{1}
For all binary words $a$, $b$, $u$, the pair encoding satisfies
$\mathrm{pair}(a \mathbin{+\!\!+} b, u) = \overline{a} \mathbin{+\!\!+} \mathrm{pair}(b, u)$,
where $\overline{a}$ is the word obtained from $a$ by doubling every bit.
\end{lemma}

\begin{lemma}[Length of a doubled word]
\lean{Complexity.taut_double_length}
\label{Complexity.taut_double_length}
\leanok
\uses{Turing.pairEncode, Turing.universal_pair_length}
\difficulty{3}
For every binary word $a$, the word obtained from $a$ by replacing each bit $v$ with
$v\,v$ has length exactly $2|a|$.
\end{lemma}

\begin{lemma}[Unary index walk]
\lean{Complexity.taut_unary_run}
\label{Complexity.taut_unary_run}
\leanok
\uses{Complexity.tautCfg, Complexity.tautTM, Complexity.taut_eval_double, Turing.MultiTapeTM.runFrom, Turing.MultiTapeTM.runFrom_add, Turing.pairEncode, Turing.universal_pair_length}
\difficulty{4}
Let $k \in \mathbb{N}$, let $r, u, p$ be binary words, and let the input be
$x = p \mathbin{+\!\!+} \mathrm{pair}(1^k \mathbin{+\!\!+} r, u)$. Started in the
unary-index evaluation state with accumulator $c$, $|p|$ input bits consumed, work word
$u$, and work head $h$, the non-tautology verifier after exactly $2k$ steps is in the
same state with $|p| + 2k$ bits consumed and work head $h + k$.
\end{lemma}

\begin{lemma}[Literal evaluation run]
\lean{Complexity.taut_literal_run}
\label{Complexity.taut_literal_run}
\leanok
\uses{Complexity.tautAssignment, Complexity.tautCfg, Complexity.tautEvalAction, Complexity.tautTM, Complexity.taut_double_length, Complexity.taut_eval_double, Complexity.taut_pair_append, Complexity.taut_unary_run, Complexity.taut_work_rewind, Std.Sat.CNF.serializeLit, Turing.FinTM.bufferTape_nat, Turing.MultiTapeTM.runFrom, Turing.MultiTapeTM.runFrom_add, Turing.pairEncode, Turing.universal_pair_length}
\difficulty{6}
Let $\ell = (v, b)$ be a literal and $u$ a word with $v < |u|$, and let the input be
$x = p \mathbin{+\!\!+} \mathrm{pair}(\sigma \mathbin{+\!\!+} r, u)$ where $\sigma$ is the
serialization of $\ell$. Started in the clause-level evaluation state with accumulator
$c$, $|p|$ bits consumed, work word $u$, and head $0$, the non-tautology verifier after
exactly $3v + 7$ steps is in the clause-level state with accumulator
$c \wedge (\alpha_u(v) = b)$, where $\alpha_u$ reads $v$ off the $v$-th bit of $u$, with
$|p| + 2|\sigma|$ bits consumed and head $0$.
\end{lemma}

\begin{lemma}[Term evaluation run]
\lean{Complexity.taut_clause_run}
\label{Complexity.taut_clause_run}
\leanok
\uses{Complexity.tautAssignment, Complexity.tautCfg, Complexity.tautTM, Complexity.taut_double_length, Complexity.taut_eval_double, Complexity.taut_literal_run, Complexity.taut_pair_append, Std.Sat.CNF.serializeClause, Std.Sat.CNF.serializeLit, Turing.MultiTapeTM.runFrom, Turing.MultiTapeTM.runFrom_add, Turing.pairEncode, Turing.universal_pair_length}
\difficulty{5}
Let $C$ be a list of literals whose variable indices are all below $|u|$ for a word
$u$, let $\sigma$ be its serialization, and let the input be
$x = p \mathbin{+\!\!+} \mathrm{pair}(\sigma \mathbin{+\!\!+} r, u)$. Started in the
clause-level evaluation state with accumulator $c$, $|p|$ bits consumed, work word $u$,
and head $0$, the non-tautology verifier within some $t \le 3|\sigma|$ steps has
consumed $|p| + 2|\sigma|$ bits, has head $0$, and is in the verdict state with bit $0$
if $c$ holds and every literal of $C$ is satisfied by the assignment read off $u$, and
in the formula-level evaluation state otherwise.
\end{lemma}

\begin{lemma}[Formula evaluation run]
\lean{Complexity.taut_formula_run}
\label{Complexity.taut_formula_run}
\leanok
\uses{Complexity.tautAssignment, Complexity.tautCfg, Complexity.tautTM, Complexity.taut_clause_run, Complexity.taut_double_length, Complexity.taut_eval_double, Complexity.taut_pair_append, Complexity.taut_verdict, Std.Sat.CNF.serialize, Std.Sat.CNF.serializeClause, Std.Sat.DNF.eval, Turing.MultiTapeTM.runFrom, Turing.MultiTapeTM.runFrom_add, Turing.MultiTapeTM.runFrom_succ_eq_step', Turing.pairEncode, Turing.universal_pair_length}
\difficulty{5}
Let $\varphi$ be a formula whose variable indices are all below $|u|$ for a word $u$,
let $\sigma$ be its serialization, and let the input be
$x = p \mathbin{+\!\!+} \mathrm{pair}(\sigma, u)$. Started in the formula-level
evaluation state with $|p|$ bits consumed, work word $u$, and head $0$, the
non-tautology verifier halts within some $t \le 3|\sigma| + 1$ steps, with work head
$0$ and output the single bit that negates the truth value of $\varphi$ read as a DNF
under the assignment $\alpha_u$ reading variable $v$ off the $v$-th bit of $u$ (false
beyond $|u|$).
\end{lemma}

\begin{lemma}[Members are bounded by the maximum]
\lean{Complexity.taut_le_max}
\label{Complexity.taut_le_max}
\leanok
\difficulty{3}
For every list $s$ of natural numbers and every member $n$ of $s$, $n$ is at most the
maximum of $s$, computed as the right fold of $\max$ with initial value $0$.
\end{lemma}

\begin{lemma}[Literal indices lie below the variable bound]
\lean{Complexity.taut_literal_lt}
\label{Complexity.taut_literal_lt}
\leanok
\uses{Complexity.taut_le_max, Std.Sat.CNF.numVars}
\difficulty{3}
Let $\varphi$ be a CNF formula, $C$ a clause of $\varphi$, and $\ell = (v, b)$ a literal
of $C$. Then $v$ is strictly less than the variable bound of $\varphi$, i.e.\ one more
than the largest variable index occurring in $\varphi$.
\end{lemma}

\begin{lemma}[Verifier correctness on correctly sized pairs]
\lean{Complexity.taut_machine_pair}
\label{Complexity.taut_machine_pair}
\leanok
\uses{Complexity.tautAssignment, Complexity.tautTM, Complexity.taut_evaluation_start, Complexity.taut_formula_run, Complexity.taut_literal_lt, Complexity.taut_machine_malformed, Complexity.taut_parse_shape, Std.Sat.CNF.decode, Std.Sat.CNF.decode_serialize, Std.Sat.CNF.numVars_decode_le, Std.Sat.CNF.parse, Std.Sat.CNF.parse_serialize, Std.Sat.CNF.serialize, Std.Sat.DNF.decode, Std.Sat.DNF.decode_cnf_serialize, Std.Sat.DNF.eval, Turing.FinTM.ComputesInTime, Turing.FinTM.ComputesInTime.mono, Turing.FinTM.computesInTime_iff, Turing.MultiTapeTM.runFrom_add, Turing.pairEncode, Turing.universal_pair_length}
\difficulty{5}
Let $x$ and $u$ be binary words with $|u| = |x| + 1$. The non-tautology verifier, run on
$\mathrm{pair}(x, u)$, halts within $10(|\mathrm{pair}(x,u)| + 1)$ steps with output
the single bit that negates the truth value of the DNF decoded from $x$ under the
assignment $\alpha_u$ reading variable $v$ off the $v$-th bit of $u$ (false beyond
$|u|$).
\end{lemma}

\begin{definition}[Split of a verifier input]
\lean{Complexity.tautSplit}
\label{Complexity.tautSplit}
\leanok
\uses{Turing.pairEncode, Turing.solveSplit}
Given a binary string $z$, its split is the pair encoding of the first $i$ bits of $z$
with the remaining $i + 1$ bits when $|z| = i + (i + 1)$ for some $i$, and the empty
word when no such $i$ exists.
\end{definition}

\begin{lemma}[The verifier handles every split]
\lean{Complexity.taut_machine_split}
\label{Complexity.taut_machine_split}
\leanok
\uses{Complexity.tautSplit, Complexity.tautTM, Complexity.tautVerifierBit, Complexity.taut_machine_empty, Complexity.taut_machine_pair, Complexity.taut_split_some, Turing.FinTM.ComputesInTime, Turing.FinTM.ComputesInTime.mono, Turing.solveSplit, Turing.universal_pair_length}
\difficulty{4}
For every binary string $z$, let $z'$ be the pair encoding of the first $i$ bits of $z$
with the rest if $|z| = 2i + 1$, and the empty word if $|z|$ is even. The
non-tautology verifier, run on $z'$, halts within $40(|z| + 1)$ steps with output the
single bit that is $0$ if $|z|$ is even and otherwise negates the truth value of the
DNF decoded from the first $i$ bits of $z$ under the assignment read off the remaining
bits.
\end{lemma}

\begin{lemma}[Composition with a machine correct on an image]
\lean{Complexity.taut_comp_on_image}
\label{Complexity.taut_comp_on_image}
\leanok
\uses{Turing.Cfg.init, Turing.FinTM, Turing.FinTM.ComputesFunInTime, Turing.FinTM.ComputesInTime, Turing.FinTM.ComputesInTime.mono, Turing.FinTM.VirtualTag, Turing.FinTM.bufferedCompTM, Turing.FinTM.bufferedComp_start, Turing.FinTM.bufferedSecondCfg, Turing.FinTM.bufferedSecondCfg_run, Turing.FinTM.computesInTime_iff, Turing.MultiTapeTM.initCfg, Turing.MultiTapeTM.output_length_le, Turing.MultiTapeTM.runFrom_add}
\difficulty{5}
Let $M$ and $U$ be finite multi-tape Turing machines over the binary alphabet, let
$f, g$ be functions on binary strings, and let $T_1, T_2 : \mathbb{N} \to \mathbb{N}$.
If $M$ computes $f$ within $T_1$, and for every input $x$ the machine $U$ run on $f(x)$
halts within $T_2(|x|)$ steps with output $g(x)$, then some finite machine computes $g$
within time $n \mapsto 2\,T_1(n) + T_2(n) + 2$.
\end{lemma}

\begin{lemma}[The non-tautology verifier language is in P]
\lean{Complexity.taut_verifier_mem_P}
\label{Complexity.taut_verifier_mem_P}
\leanok
\uses{Complexity.P, Complexity.mem_P_iff, Complexity.tautSplit, Complexity.tautTM, Complexity.tautVerifier, Complexity.tautVerifierBit, Complexity.taut_comp_on_image, Complexity.taut_machine_split, Turing.FinTM.ComputesFunInTime, Turing.FinTM.ComputesInTime, Turing.FinTM.ComputesInTime.mono, Turing.FinTM.computesFunInTime_splitSolve, Turing.MultiTapeTM.indicator, Turing.solveSplit}
\difficulty{5}
The set of binary strings $z$ such that $|z| = i + (i + 1)$ for some $i$ and the DNF
decoded from the first $i$ bits of $z$ evaluates to false under the assignment reading
variable $v$ off the $v$-th of the remaining $i + 1$ bits (false beyond them) belongs
to $\mathrm{P}$.
\end{lemma}

\begin{theorem}[TAUTOLOGY is in coNP]
\lean{Complexity.TAUTOLOGY_mem_coNP}
\label{Complexity.TAUTOLOGY_mem_coNP}
\leanok
\uses{Complexity.TAUTOLOGY, Complexity.coNP, Complexity.taut_membership_of_verifier, Complexity.taut_verifier_mem_P}
\difficulty{1}
The language $\mathrm{TAUTOLOGY}$, of binary strings whose decoded DNF formula is true
under every assignment, belongs to coNP; that is, its complement is in
$\mathrm{NP}$ [AB09, Section~2.6.1].
\end{theorem}

\begin{definition}[Polarity-flipping output bit]
\lean{Complexity.tautDualBit}
\label{Complexity.tautDualBit}
\leanok
\uses{Complexity.TautSyntax}
Given a state $q$ of the six-state formula syntax scanner and a bit $b$, the
polarity-flipping output bit is $\neg b$ when $q$ is the polarity state and $b$
otherwise.
\end{definition}

\begin{definition}[Polarity-flipping scan]
\lean{Complexity.tautDualScan}
\label{Complexity.tautDualScan}
\leanok
\uses{Complexity.TautSyntax, Complexity.tautDualBit, Complexity.tautSyntaxStep}
Given a state $q$ of the six-state formula syntax scanner and a binary string $x$, the
polarity-flipping scan is the string obtained by running the scanner over $x$ from $q$
and outputting each bit as read, except that a bit read in the polarity state is
flipped.
\end{definition}

\begin{lemma}[Polarity-flipping scan of a unary run]
\lean{Complexity.tautDual_unary}
\label{Complexity.tautDual_unary}
\leanok
\uses{Complexity.tautDualBit, Complexity.tautDualScan, Complexity.tautSyntaxStep}
\difficulty{3}
For every $n \in \mathbb{N}$ and binary string $r$, the polarity-flipping scan from the
unary-index state of the formula syntax scanner maps $1^n \mathbin{+\!\!+} r$ to
$1^n$ followed by the polarity-flipping scan of $r$ from the unary-index state.
\end{lemma}

\begin{lemma}[Polarity-flipping scan of a literal]
\lean{Complexity.tautDual_literal}
\label{Complexity.tautDual_literal}
\leanok
\uses{Complexity.tautDualBit, Complexity.tautDualScan, Complexity.tautDual_unary, Complexity.tautSyntaxStep, Std.Sat.CNF.serializeLit}
\difficulty{3}
For every literal $(v, b)$ and binary string $r$, the polarity-flipping scan from the
clause-level state of the formula syntax scanner maps the serialization of $(v, b)$
followed by $r$ to the serialization of $(v, \neg b)$ followed by the
polarity-flipping scan of $r$ from the clause-level state.
\end{lemma}

\begin{lemma}[Polarity-flipping scan of a clause]
\lean{Complexity.tautDual_clause}
\label{Complexity.tautDual_clause}
\leanok
\uses{Complexity.tautDualScan, Complexity.tautDual_literal, Std.Sat.CNF.serializeClause, Std.Sat.CNF.serializeLit}
\difficulty{4}
For every clause $C$ and binary string $r$, the polarity-flipping scan from the
clause-level state of the formula syntax scanner maps the serialization of $C$
followed by $r$ to the serialization of the clause obtained from $C$ by negating every
literal's polarity, followed by the polarity-flipping scan of $r$ from the
formula-level state.
\end{lemma}

\begin{lemma}[Polarity-flipping scan of a serialized formula]
\lean{Complexity.tautDual_serialize}
\label{Complexity.tautDual_serialize}
\leanok
\uses{Complexity.tautDualScan, Complexity.tautDual_clause, Std.Sat.CNF.dual, Std.Sat.CNF.serialize, Std.Sat.CNF.serializeClause, Std.Sat.DNF.serialize}
\difficulty{4}
For every CNF formula $\varphi$, the polarity-flipping scan from the formula-level
state of the formula syntax scanner maps the serialization of $\varphi$ to the
serialization of its De Morgan dual, the DNF obtained from $\varphi$ by negating every
literal in place.
\end{lemma}

\begin{definition}[Guarded dual output]
\lean{Complexity.tautDualOutput}
\label{Complexity.tautDualOutput}
\leanok
\uses{Complexity.tautDualScan, Complexity.tautSyntaxStep}
Given a binary string $x$, the guarded dual output is the polarity-flipping scan of $x$
from the formula-level state when the six-state formula syntax scan of $x$ from
formula level ends in the exact-end state, and the one-bit word $0$ (the serialization
of the empty formula) otherwise.
\end{definition}

\begin{lemma}[Guarded output is decode--dual--serialize]
\lean{Complexity.tautDual_output}
\label{Complexity.tautDual_output}
\leanok
\uses{Complexity.tautDualOutput, Complexity.tautDual_serialize, Complexity.tautSyntaxAccept, Complexity.tautSyntaxStep, Complexity.taut_parse_shape, Complexity.taut_syntax_parse, Std.Sat.CNF.decode, Std.Sat.CNF.dual, Std.Sat.CNF.parse, Std.Sat.DNF.serialize}
\difficulty{4}
For every binary string $x$, the guarded dual output of $x$ — the polarity-flipping
scan of $x$ if $x$ passes the formula syntax scan, and the word $0$ otherwise — equals
the serialization of the De Morgan dual (every literal negated, read as a DNF) of the
CNF formula decoded from $x$.
\end{lemma}

\begin{definition}[Control states of the dual transducer]
\lean{Complexity.TautDualControl}
\label{Complexity.TautDualControl}
\leanok
\uses{Complexity.TautSyntax}
The control states of the dual transducer are an anchor state, scanning states and
copying states each carrying a formula-scanner state, and rewind states carrying a
flag that selects whether the rewind ends at the anchor or starts the copying pass.
\end{definition}

\begin{definition}[Transducer action]
\lean{Complexity.tautDualAction}
\label{Complexity.tautDualAction}
\leanok
\uses{Complexity.TautDualControl, Turing.Action}
Given an input-head displacement $d \in \{-1, 0, 1\}$, a control state $q$, and an
optional output bit $o$ (absent by default), the transducer action is the action of a
machine with no work tapes that moves the input head by $d$, emits $o$ if present, and
enters $q$.
\end{definition}

\begin{definition}[De Morgan dual transducer machine]
\lean{Complexity.tautDualTM}
\label{Complexity.tautDualTM}
\leanok
\uses{Complexity.TautDualControl, Complexity.tautDualAction, Complexity.tautDualBit, Complexity.tautSyntaxStep, Turing.FinTM}
The dual transducer is a finite Turing machine over the binary alphabet with no work
tapes. From its anchor state it scans the whole input silently with the six-state
formula syntax scanner. If the scan ends in the exact-end state, it rewinds and makes
a second pass emitting every input bit, flipped exactly when read in the polarity
state; otherwise it emits the single bit $0$. In both cases it then rewinds to the
first input position and re-enters the anchor state.
\end{definition}

\begin{definition}[Dual transducer configuration]
\lean{Complexity.tautDualCfg}
\label{Complexity.tautDualCfg}
\leanok
\uses{Complexity.TautDualControl, Turing.Cfg}
Given an input $x$, a control state $q$, an input-head position
$p \in \{0, \dots, |x| + 1\}$ (positions $0$ and $|x|+1$ being the boundary blanks),
and an output word $o$, the transducer configuration is the configuration of the
zero-work-tape dual transducer on $x$ with state $q$, input head at $p$, and output
$o$.
\end{definition}

\begin{lemma}[Effect of a transducer action]
\lean{Complexity.tautDual_action}
\label{Complexity.tautDual_action}
\leanok
\uses{Complexity.TautDualControl, Complexity.tautDualAction, Complexity.tautDualCfg, Turing.Action.apply, Turing.Cfg.ext_zero_tapes, Turing.moveInputPos}
\difficulty{1}
For every input $x$, applying the zero-work-tape action that moves the input head by
$d$, enters $q'$, and optionally emits $o$ to the configuration with state $q$, head
position $p$, and output $w$ yields the configuration with state $q'$, head position
$p$ moved by $d$ (clamped at the boundary blanks), and output $w$ followed by $o$.
\end{lemma}

\begin{lemma}[Input symbol of a transducer configuration]
\lean{Complexity.tautDual_read}
\label{Complexity.tautDual_read}
\leanok
\uses{Complexity.TautDualControl, Complexity.tautDualCfg, Turing.Cfg.inputSymbol, Turing.FinTM.inputSymbol_at}
\difficulty{1}
Let $x$ be an input and $i \le |x|$. In every configuration of the dual transducer on
$x$ with input head at position $i + 1$, the symbol read is the $i$-th bit of $x$ when
$i < |x|$ and the right boundary blank when $i = |x|$.
\end{lemma}

\begin{definition}[Anchor-free run segment]
\lean{Complexity.tautDualSegment}
\label{Complexity.tautDualSegment}
\leanok
\uses{Complexity.TautDualControl, Complexity.tautDualTM, Turing.Cfg, Turing.MultiTapeTM.runFrom}
Given configurations $c$ and $d$ of the dual transducer on an input $x$ and a time
$t \in \mathbb{N}$, the segment predicate holds when running the transducer for $t$
steps from $c$ reaches exactly $d$, and for every $j < t$ the configuration after $j$
steps from $c$ is not in the anchor state.
\end{definition}

\begin{lemma}[Empty segment]
\lean{Complexity.tautDual_segment_zero}
\label{Complexity.tautDual_segment_zero}
\leanok
\uses{Complexity.TautDualControl, Complexity.tautDualSegment, Turing.Cfg}
\difficulty{1}
For every configuration $c$ of the dual transducer, running it for $0$ steps from $c$
reaches $c$ without visiting the anchor state strictly before, so $c$ to $c$ is a
segment of length $0$.
\end{lemma}

\begin{lemma}[Single-step segment]
\lean{Complexity.tautDual_segment_one}
\label{Complexity.tautDual_segment_one}
\leanok
\uses{Complexity.TautDualControl, Complexity.tautDualSegment, Complexity.tautDualTM, Turing.Cfg, Turing.MultiTapeTM.step}
\difficulty{2}
Let $c$ and $d$ be configurations of the dual transducer such that one step takes $c$
to $d$ and $c$ is not in the anchor state. Then one step from $c$ reaches $d$ with no
anchor visit before it, i.e.\ $c$ to $d$ is an anchor-free segment of length $1$.
\end{lemma}

\begin{lemma}[Segments compose]
\lean{Complexity.tautDual_segment_add}
\label{Complexity.tautDual_segment_add}
\leanok
\uses{Complexity.TautDualControl, Complexity.tautDualSegment, Turing.Cfg, Turing.MultiTapeTM.runFrom_add}
\difficulty{4}
Let $c, d, e$ be configurations of the dual transducer. If running $a$ steps from $c$
reaches $d$ without visiting the anchor state before time $a$, and running $b$ steps
from $d$ reaches $e$ without visiting the anchor before time $b$, then running $a + b$
steps from $c$ reaches $e$ without visiting the anchor before time $a + b$.
\end{lemma}

\begin{lemma}[Silent syntax scan of a suffix]
\lean{Complexity.tautDual_scan}
\label{Complexity.tautDual_scan}
\leanok
\uses{Complexity.TautSyntax, Complexity.tautDualAction, Complexity.tautDualCfg, Complexity.tautDualSegment, Complexity.tautDualTM, Complexity.tautDual_action, Complexity.tautDual_read, Complexity.tautDual_segment_add, Complexity.tautDual_segment_one, Complexity.tautDual_segment_zero, Complexity.tautSyntaxStep, Turing.Action.apply, Turing.MultiTapeTM.step, Turing.moveInputPos_pos_of_ne_right}
\difficulty{5}
Let the input be $x = p \mathbin{+\!\!+} r$, let $q$ be a formula-scanner state, and
let $o$ be an output word. Started in the scanning state carrying $q$ at head position
$|p| + 1$ with output $o$, the dual transducer reaches the scanning state carrying the
fold of the scanner transition over $r$ from $q$, at position $|x| + 1$ with output
still $o$, in exactly $|r|$ steps and without visiting the anchor state before then.
\end{lemma}

\begin{lemma}[Emitting copy pass over a suffix]
\lean{Complexity.tautDual_copy}
\label{Complexity.tautDual_copy}
\leanok
\uses{Complexity.TautSyntax, Complexity.tautDualAction, Complexity.tautDualBit, Complexity.tautDualCfg, Complexity.tautDualScan, Complexity.tautDualSegment, Complexity.tautDualTM, Complexity.tautDual_action, Complexity.tautDual_read, Complexity.tautDual_segment_add, Complexity.tautDual_segment_one, Complexity.tautDual_segment_zero, Complexity.tautSyntaxStep, Turing.Action.apply, Turing.MultiTapeTM.step, Turing.moveInputPos_pos_of_ne_right}
\difficulty{5}
Let the input be $x = p \mathbin{+\!\!+} r$, let $q$ be a formula-scanner state, and
let $o$ be an output word. Started in the copying state carrying $q$ at head position
$|p| + 1$ with output $o$, the dual transducer reaches, in exactly $|r|$ steps and
without visiting the anchor state before then, the copying state carrying the fold of
the scanner transition over $r$ from $q$, at position $|x| + 1$, with output $o$
followed by the polarity-flipping scan of $r$ from $q$.
\end{lemma}

\begin{lemma}[Silent rewind to the start]
\lean{Complexity.tautDual_rewind}
\label{Complexity.tautDual_rewind}
\leanok
\uses{Complexity.tautDualAction, Complexity.tautDualCfg, Complexity.tautDualSegment, Complexity.tautDualTM, Complexity.tautDual_action, Complexity.tautDual_read, Complexity.tautDual_segment_add, Complexity.tautDual_segment_one, Turing.Action.apply, Turing.Cfg.inputSymbol, Turing.MultiTapeTM.step, Turing.moveInputPos_neg_of_ne_left, Turing.moveInputPos_pos_of_ne_right}
\difficulty{5}
Let $x$ be an input, $o$ an output word, $f$ a flag, and $p \le |x|$. Started in the
rewind state with flag $f$ at head position $p$ and output $o$, the dual transducer
reaches head position $1$ with output $o$, in the anchor state if $f$ is true and in
the formula-level copying state otherwise, in exactly $p + 1$ steps and without
visiting the anchor before then.
\end{lemma}

\begin{lemma}[End of the copying pass]
\lean{Complexity.tautDual_copy_end}
\label{Complexity.tautDual_copy_end}
\leanok
\uses{Complexity.TautSyntax, Complexity.tautDualAction, Complexity.tautDualCfg, Complexity.tautDualSegment, Complexity.tautDualTM, Complexity.tautDual_action, Complexity.tautDual_read, Complexity.tautDual_segment_one, Turing.Action.apply, Turing.moveInputPos_neg_of_ne_left}
\difficulty{3}
For every input $x$, output word $o$, and formula-scanner state $q$, the configuration
reached by one step of the dual transducer from the copying state carrying $q$ at the
right boundary position $|x| + 1$ with output $o$ is the rewind state with flag true
at position $|x|$ with output still $o$, and the anchor state is not visited.
\end{lemma}

\begin{lemma}[Round on a well-formed input]
\lean{Complexity.tautDual_valid}
\label{Complexity.tautDual_valid}
\leanok
\uses{Complexity.tautDualAction, Complexity.tautDualCfg, Complexity.tautDualScan, Complexity.tautDualSegment, Complexity.tautDualTM, Complexity.tautDual_action, Complexity.tautDual_copy, Complexity.tautDual_copy_end, Complexity.tautDual_read, Complexity.tautDual_rewind, Complexity.tautDual_scan, Complexity.tautDual_segment_add, Complexity.tautDual_segment_one, Complexity.tautSyntaxStep, Turing.Action.apply, Turing.moveInputPos_neg_of_ne_left}
\difficulty{5}
Let $x$ be a binary string on which the six-state formula syntax scan from formula
level ends in the exact-end state. Started in the formula-level scanning state at head
position $1$ with empty output, the dual transducer reaches the anchor state at
position $1$ with output the polarity-flipping scan of $x$, in exactly $4|x| + 4$ steps
and without visiting the anchor before then.
\end{lemma}

\begin{lemma}[Round on a malformed input]
\lean{Complexity.tautDual_invalid}
\label{Complexity.tautDual_invalid}
\leanok
\uses{Complexity.tautDualAction, Complexity.tautDualCfg, Complexity.tautDualSegment, Complexity.tautDualTM, Complexity.tautDual_action, Complexity.tautDual_read, Complexity.tautDual_rewind, Complexity.tautDual_scan, Complexity.tautDual_segment_add, Complexity.tautDual_segment_one, Complexity.tautSyntaxStep, Turing.Action.apply, Turing.moveInputPos_neg_of_ne_left}
\difficulty{4}
Let $x$ be a binary string on which the six-state formula syntax scan from formula
level does not end in the exact-end state. Started in the formula-level scanning state
at head position $1$ with empty output, the dual transducer reaches the anchor state
at position $1$ with output the single bit $0$, in exactly $2|x| + 2$ steps and
without visiting the anchor before then.
\end{lemma}

\begin{lemma}[One complete transducer round]
\lean{Complexity.tautDual_round}
\label{Complexity.tautDual_round}
\leanok
\uses{Complexity.tautDualCfg, Complexity.tautDualOutput, Complexity.tautDualSegment, Complexity.tautDualTM, Complexity.tautDual_invalid, Complexity.tautDual_valid, Complexity.tautSyntaxStep, Turing.Cfg.ext_zero_tapes, Turing.Cfg.ofWords, Turing.MultiTapeTM.runFrom, Turing.MultiTapeTM.runFrom_succ_eq_step, Turing.MultiTapeTM.step, Turing.moveInputPos_zero, Turing.stateWord}
\difficulty{5}
For every binary string $x$, let $c$ be the dual transducer's configuration on $x$ in
the anchor state with input head at position $1$ and empty output. There is a time $t$
with $0 < t \le 4|x| + 5$ such that the configurations at all times $j$ with
$0 < j < t$ are not in the anchor state, and the configuration at time $t$ equals $c$
except that its output is the guarded dual output of $x$: the polarity-flipping scan
of $x$ if $x$ passes the formula syntax scan, and the single bit $0$ otherwise.
\end{lemma}

\begin{lemma}[Decode--dual--serialize is polynomial-time]
\lean{Complexity.tautDual_poly}
\label{Complexity.tautDual_poly}
\leanok
\uses{Complexity.PolyTimeComputable, Complexity.tautDualOutput, Complexity.tautDualTM, Complexity.tautDual_output, Complexity.tautDual_round, Std.Sat.CNF.decode, Std.Sat.CNF.dual, Std.Sat.DNF.serialize, Turing.Cfg.ext_zero_tapes, Turing.FinTM.ComputesFunInTime, Turing.FinTM.ComputesInTime, Turing.FinTM.ComputesInTime.mono, Turing.FinTM.computesFunInTime_const, Turing.FinTM.exists_emitLoopTM}
\difficulty{5}
The function that maps a binary string $x$ to the serialization of the De Morgan dual
of the CNF formula decoded from $x$ — the DNF obtained by negating every literal in
place — is polynomial-time computable.
\end{lemma}

\begin{theorem}[TAUTOLOGY is coNP-complete]
\lean{Complexity.TAUTOLOGY_coNPComplete}
\label{Complexity.TAUTOLOGY_coNPComplete}
\leanok
\uses{Complexity.PolyTimeComputable.comp, Complexity.SAT_NPHard, Complexity.TAUTOLOGY, Complexity.TAUTOLOGY_mem_coNP, Complexity.coNPComplete, Complexity.tautDual_poly, Std.Sat.CNF.Satisfiable, Std.Sat.CNF.decode, Std.Sat.CNF.dual, Std.Sat.CNF.tautology_dual_iff, Std.Sat.DNF.Tautology, Std.Sat.DNF.decode, Std.Sat.DNF.decode_serialize, Std.Sat.DNF.serialize}
\difficulty{4}
The language $\mathrm{TAUTOLOGY}$, of binary strings whose decoded DNF formula is true
under every assignment, is coNP-complete: it belongs to coNP, and
every language in coNP is polynomial-time Karp reducible to it
[AB09, Example~2.21].
\end{theorem}
